% ===================================================================
\section{Donor-level battery tests}\label{app:breaks}
% ===================================================================
The qualitative audit of Section~\ref{sec:donoraudit} is the primary identification defence.
This appendix reports its empirical counterpart, a battery of donor-level tests run on each
candidate's own return or level series before any model is fit. The battery asks two
questions. The first is whether a donor jumps at the event, a direct reading of the
no-interference assumption. The second is whether the donor's relationship with Brent stays
stable across the fitting window, since an unstable loading would distort the synthetic.
Every test is model-agnostic and works one donor at a time, so it confirms or supplements the
audit rather than replacing it. Throughout we write $r_{j,t}$ for donor $j$'s daily log-return,
$y_t$ for Brent's, $T_0$ for the known cutoff, and $m=32$ for the number of candidates.

\subsection{Interference at the event}
These tests read each donor's own returns for a shift at $T_0$, with $n_{\text{pre}}$ and
$n_{\text{post}}$ the counts of pre-cutoff and post-cutoff observations. Two tests carry the
flag and a third is reported alongside.

\begin{enumerate}
\item \textbf{Permutation mean-shift test.} The statistic is the post-minus-pre difference in
mean return,
\begin{equation}
d_j=\frac{1}{n_{\text{post}}}\sum_{t>T_0} r_{j,t}\;-\;\frac{1}{n_{\text{pre}}}\sum_{t\le T_0} r_{j,t}.
\label{eq:permshift}
\end{equation}
Its null distribution comes from reshuffling the pre and post labels $B$ times and recomputing
the statistic on each draw, which gives the p-value
\begin{equation}
p_j=\frac{1}{B}\sum_{b=1}^{B}\mathbf{1}\big\{\,|d_j^{(b)}|\ge|d_j|\,\big\}.
\label{eq:permp}
\end{equation}
This is the known-date permutation counterpart of the structural-break test of
\citet{andrews1993} and needs no normality assumption.

\item \textbf{Wilcoxon event-window test.} Let $E$ be a short window of returns around $T_0$ and
$C$ a sixty-day pre-window control, with sizes $n_E$ and $n_C$ and pooled size $N=n_E+n_C$.
Ranking the pooled returns and summing the ranks that fall in $E$ gives the rank-sum
\begin{equation}
W=\sum_{t\in E}\operatorname{rank}(r_{j,t}),\qquad \mathbb{E}_0[W]=\tfrac{1}{2}\,n_E(N+1),
\label{eq:wilcoxon}
\end{equation}
and the two-sided p-value follows from the Mann--Whitney null. Working in ranks makes it robust
to the heavy tails of daily returns.

\item \textbf{Benjamini--Hochberg correction.} Each test returns one p-value per donor across the
$m$ candidates. Ordering them $p_{(1)}\le\dots\le p_{(m)}$, the procedure rejects every hypothesis
up to
\begin{equation}
k^\star=\max\Big\{\,k:\;p_{(k)}\le\tfrac{k}{m}\,q\,\Big\},\qquad q=0.10,
\label{eq:bh}
\end{equation}
which controls the false-discovery rate at $q$ \citep{benjaminihochberg1995}. A donor is flagged
as potentially treated when either the permutation or the Wilcoxon test rejects after this
correction.

\item \textbf{Kolmogorov--Smirnov distance} (reported, not in the flag rule). The two-sample
statistic is the largest gap between the pre and post empirical distributions of the donor's
returns,
\begin{equation}
D_j=\sup_x\big|\hat F^{\text{pre}}_j(x)-\hat F^{\text{post}}_j(x)\big|.
\label{eq:ks}
\end{equation}
At these sample sizes it rejects on any broad macro-regime shift rather than on event-treatment
of the individual donor, so we report it but keep it out of the flag rule.
\end{enumerate}

For Hormuz 2026 the battery agrees with the audit on 31 of the 32 candidates, with no flag
surviving the correction. The single disagreement is Cotton, which the audit drops on the
indirect oil-to-polyester channel that the contemporaneous tests are not built to reproduce.
For Russia 2022 the tests miss the audit's directly exposed donors, because at about 415
pre-period observations the war-driven moves are not extreme enough to survive the correction
across the candidate set, and they add one flag on CNY that traces to the concurrent China
zero-COVID dynamics rather than to the invasion. The Kolmogorov--Smirnov test flags roughly
21 of the 32 Russia donors, far more than the audit, which is the over-rejection on the 2022
regime shift that justifies keeping it out of the flag rule.

\subsection{Stability of the donor-Brent relationship}
Three tests check that the donor-Brent relationship holds across the pre-window.

\begin{enumerate}
\item \textbf{Distance-correlation change.} The distance correlation \citep{szekely2007} measures
general dependence and is zero only under independence. From the pairwise distances
$a_{kl}=|r_{j,k}-r_{j,l}|$ and $b_{kl}=|y_k-y_l|$, double-centred into $A_{kl}$ and $B_{kl}$, it is
\begin{equation}
\operatorname{dCor}(r_j,y)=\frac{\operatorname{dCov}(r_j,y)}{\sqrt{\operatorname{dVar}(r_j)\,\operatorname{dVar}(y)}},
\qquad
\operatorname{dCov}^2(r_j,y)=\frac{1}{n^2}\sum_{k,l}A_{kl}B_{kl}.
\label{eq:dcor}
\end{equation}
Splitting the pre-window into thirds $I_1,I_2,I_3$ and computing the distance correlation on each,
the statistic is the largest change across them, $\max_{a<b}|\operatorname{dCor}_a-\operatorname{dCor}_b|$.
We read it as a narrative flag rather than an exclusion.

\item \textbf{Bai--Perron break dating.} Fitting the donor-Brent return relationship
$r_{j,t}=\alpha_i+\beta_i y_t+u_t$ with $m$ unknown breaks that partition the sample, the break
dates minimise the global segment sum of squares,
\begin{equation}
(\hat T_1,\dots,\hat T_m)=\arg\min_{T_1<\dots<T_m}\sum_{i=1}^{m+1}\sum_{t=T_{i-1}+1}^{T_i}
\big(r_{j,t}-\alpha_i-\beta_i y_t\big)^2,
\label{eq:baiperron}
\end{equation}
with the number of breaks chosen by the Bayesian information criterion (BIC) \citep{baiperron1998,baiperron2003}.

\item \textbf{Harvey--Leybourne--Taylor trend-break test.} Because the relationship above
regresses on Brent, a break dated at the cutoff is ambiguous, since Brent is treated there too.
This test instead reads the donor's own log-level trend and is robust to whether the series is
stationary or integrated. It combines a statistic $t_0$ that is optimal under stationarity with a
statistic $t_1$ that is optimal under a unit root,
\begin{equation}
\mathcal{T}_\lambda=(1-\lambda)\,t_0+\lambda\,t_1,
\label{eq:hlt}
\end{equation}
where the weight $\lambda$ is driven by a KPSS stationarity statistic so the test keeps the same
asymptotic size in either case \citep{harvey2009}. The KPSS input is
\begin{equation}
\eta_j=\frac{1}{T^2\hat\omega^2}\sum_{t=1}^{T}S_t^2,\qquad S_t=\sum_{i=1}^{t}\hat e_i,
\label{eq:kpss}
\end{equation}
the scaled partial sums of the residuals $\hat e$ from regressing the donor level on a constant
and trend, with $\hat\omega^2$ a long-run variance estimate \citep{kpss1992}. The exact weight
function is given in \citet{harvey2009}.
\end{enumerate}

Most donor breaks predate both pre-windows, clustering at the post-crisis normalisation and the
2020 COVID crash, and the COVID break falls just before the Russia pre-window start, which
vindicates that start date. Inside the pre-windows the dating flags two Russia donors, Wheat,
which the audit already excludes, and the VIX on 2022-01-07 about seven weeks before the
invasion, and three Hormuz donors, Silver, Platinum, and EM equities, on a common September 2024
precious-metals move with none near the cutoff. A break inside the pre-window is informative
rather than disqualifying, so no donor is removed by this test and the VIX is retained.

For the trend-break test all 32 donors classify as integrated, so it avoids the spurious trend
breaks a naive levels regression would find, and no donor's own trend breaks within eight weeks
of either cutoff. No retained donor was itself trend-treated by either event.
